\chapter{Resumen}

La inteligencia artificial se ha consolidado como una tecnología
transformadora en sectores como la sanidad, las finanzas, la industria o la
administración pública, y la calidad de los datos que alimentan estos sistemas
constituye un factor determinante en su fiabilidad y en la equidad de sus
resultados.  Sin embargo, las herramientas disponibles para la evaluación de
dicha calidad presentan limitaciones significativas: las bibliotecas basadas
en código imponen una barrera técnica elevada, mientras que los \emph{scripts}
ad~hoc carecen de estandarización, historización y mecanismos de decisión
automatizados.

Este Trabajo Fin de Grado presenta \dataqual, una plataforma web
\emph{full-stack} para la evaluación automatizada de la calidad de datos.
La plataforma implementa siete métricas de calidad alineadas con los
estándares ISO/IEC~5259 (calidad de datos para \ac{IA}) e ISO/IEC~25012
(modelo general de calidad de datos), un motor modular y extensible, un
sistema de \emph{fingerprinting} SHA-256 para la trazabilidad de incidencias
entre evaluaciones y \emph{Quality Gates} configurables inspirados en
SonarQube.  El sistema se despliega como ocho servicios Docker orquestados
con Docker~Compose, combinando una interfaz web interactiva (Next.js, React)
con una API RESTful (Flask, Python) y procesamiento asíncrono
(Celery, Redis).

La plataforma ha sido validada con conjuntos de datos ampliamente utilizados
en el ámbito de la \ac{IA} y la ciencia de datos, confirmando la detección
efectiva de problemas de calidad, la trazabilidad de incidencias entre
evaluaciones sucesivas y la correcta clasificación de datasets mediante
\emph{Quality Gates} con estados PASSED, WARNING y FAILED.

\dataqual contribuye a democratizar el acceso a la evaluación estandarizada
de calidad de datos, eliminando la barrera técnica que suponen las
herramientas basadas en código.  Su alineación con el Reglamento Europeo de
Inteligencia Artificial facilita el cumplimiento normativo en un contexto
regulatorio cada vez más exigente, contribuyendo al desarrollo de sistemas
de \ac{IA} más fiables, equitativos y auditables.


\chapter{Abstract}

Artificial intelligence has become a transformative technology in sectors
such as healthcare, finance, industry and public administration, and the
quality of the data that feed these systems is a determining factor in their
reliability and the fairness of their outcomes.  However, the tools available
for data quality assessment have significant limitations: code-based
libraries impose a high technical barrier, while ad~hoc scripts lack
standardisation, historisation and automated decision mechanisms.

This undergraduate thesis presents \dataqual, a full-stack web platform for
automated data quality evaluation.  The platform implements seven quality
metrics aligned with the ISO/IEC~5259 (data quality for AI) and ISO/IEC~25012
(general data quality model) standards, a modular and extensible engine,
a SHA-256 fingerprinting system for issue traceability across evaluations,
and configurable Quality Gates inspired by SonarQube.  The system is deployed
as eight Docker services orchestrated with Docker~Compose, combining an
interactive web interface (Next.js, React) with a RESTful API
(Flask, Python) and asynchronous processing (Celery, Redis).

The platform has been validated with widely used datasets in the field of
AI and data science, confirming the effective detection of quality issues,
issue traceability across successive evaluations, and correct dataset
classification through Quality Gates with PASSED, WARNING and FAILED states.

\dataqual contributes to democratising access to standardised data quality
evaluation by removing the technical barrier imposed by code-based tools.
Its alignment with the European Artificial Intelligence Act supports
regulatory compliance in an increasingly demanding landscape, fostering the
development of more reliable, fair and auditable AI systems.
